A photon triggers an avalanche in a SPAD by generating an electron-hole pair responsible for causing an ionization chain reaction inside the strong depletion region of the diode. 
The avalanche spreads toward the diode extremes, where we can collect the charge carriers and detect a signal. 
Afterwards, the SPAD must reduce the voltage in the diode to stop the avalanche in a process called quenching. 
The time period between the start and the end of the avalanche is referred to as \textit{dead-time} because the detector remains blind to subsequent incident photons. 
The detection process concludes after the SPAD is recharged to its Operating Voltage (VOP).\\

Defects in the silicon lattice can trap charges which cause fake \textit{afterpulsing} clicks immediately following recharge. 
In this practical, we aim to run a correlation measurement where the start and stop channels are the same SPAD. 
The resulting histogram, known as an autocorrelation function, will map the exact temporal duration of this blind spot 
and quantify the probability of afterpulsing events by analyzing the distribution of time delays between consecutive avalanches.